%%% BEGIN fig:pdf-operator
\begin{figure}[htbp]
\centering
\tikzset{t1 hadron/.style={draw,solid,line width=1.5pt,line cap=butt}}
\begin{tikzpicture}[fd]
  \coordinate (ol) at (-1.9,2.4); \coordinate (or) at (1.9,2.4);
  \coordinate (bl) at (-1.9,0);   \coordinate (br) at (1.9,0);
  % remnant X crossing the cut (drawn first; the blobs cover the ends)
  \foreach \y in {0.24,0,-0.24} \draw[plain] (bl) ++(0,\y) -- ++(3.8,0);
  \draw[quark] (-1.9,0.43) -- (ol);
  \draw[wilson=0.25] (ol) -- (or);
  \draw[quark] (or) -- (1.9,0.43);
  \draw[t1 hadron] (-3.65,0) -- (bl);
  \draw[t1 hadron] (br) -- (3.65,0);
  \draw[hadron blob] (bl) ellipse[x radius=0.82,y radius=0.45];
  \draw[hadron blob] (br) ellipse[x radius=0.82,y radius=0.45];
  \node[opvertex] at (ol) {}; \node[opvertex] at (or) {};
  \draw[cut] (0,-0.85) -- (0,3.0);
  % labels
  \path[mom={$k$},mom pos=0.62] (-1.9,0.45) -- (ol);
  \path[mom={$k$},mom pos=0.38] (or) -- (1.9,0.45);
  \node[left,font=\small] at (-2.02,0.92) {$k^+=\xi P^+$};
  \node[above left,inner sep=3pt] at (ol) {$\psi(0)$};
  \node[above right,inner sep=3pt] at (or) {$\bar\psi(w)$};
  \node[below,font=\footnotesize,inner sep=4pt] at (-0.95,2.4) {$W[\infty n,0]$};
  \node[below,font=\footnotesize,inner sep=4pt] at (0.95,2.4) {$W[w,\infty n]$};
  \path[mom={$P$},mom sep=0.3cm] (-3.65,0) -- (-2.72,0);
  \path[mom={$P$},mom sep=0.3cm] (2.72,0) -- (3.65,0);
  \node[font=\small] at (0.42,-0.62) {$X$};
\end{tikzpicture}
\caption{夸克分布函数的割图表示。在定义
$f_{q/H}(\xi)=\int\frac{\dd w^-}{4\pi}\,\ee^{-\ii\xi P^+w^-}\langle P|\bar\psi(w)\gamma^+W[w,0]\psi(0)|P\rangle$
（$w=w^-n$，$n$ 为满足 $n\cdot v=v^+$ 的光锥矢量）中插入末态完备集并完成 $w^-$ 积分，得
$f_{q/H}(\xi)=\sum_X\tfrac12\,\delta\big(\xi P^+-(P-p_X)^+\big)\,\langle P|\bar\psi(0)W[0,\infty n]|X\rangle\,\gamma^+\,\langle X|W[\infty n,0]\psi(0)|P\rangle$。
割线左侧为振幅、右侧为共轭振幅，$X$ 为穿过割线的剩余部分。两个 $\otimes$ 为非定域算符的两端，算符顶点的 Feynman 规则为
$\tfrac12\gamma^+\delta(k^+-\xi P^+)$（配合 $\int\dd^4k/(2\pi)^4$，Wilson 线上无胶子附着时）；
双线为 Wilson 线 $W[w,0]=W[w,\infty n]\,W[\infty n,0]$，被末态割线分为振幅中的 $W[\infty n,0]$ 与共轭振幅中的 $W[w,\infty n]$。
箭头表示费米子数（色）流向。}
\label{fig:pdf-operator}
\end{figure}
%%% END fig:pdf-operator
